\chapter{Implementation: Variational Monte Carlo} \label{chp:WFE}
\epigraph{There are only two hard things in Computer Science: cache invalidation and naming things.}{Phil Karlton, \cite{fowler_bliki:_nodate}}
\iffalse
\begin{figure}[H]
	\centering
	\includegraphics[scale=0.4]{Images/example.png}
	\caption{Caption}
\end{figure}
\fi

In this chapter de will describe the implemented variational Monte Carlo (VMC) code, which was developed from scratch in C++. As the code itself is around 7000 significant\footnote{Significant lines of code in this sense means lines that are not blank or commented. Counted by the cloc program \cite{aldanial_cloc_2019}.} lines of code, we will just go through selected and often not obvious parts. As often said, \textit{good planning is half the battle}, which largely relates to writing VMC code. The code was rewritten and restructured several times before we ended on the final version. As a starting point, we used Morten Ledum's VMC framework \cite{ledum_simple_2016}, which was meant as an example implementation in the course \textit{FYS4411 - Computational Physics II: Quantum Mechanical Systems}. The entire source code can be found on the authors github, \cite{nordhagen_general_2019}.

For all matrix operations, the open source template library for linear algebra Eigen was used throughout the code. Eigen provides an elegant interface, with support for all the needed matrix and vector operations. In addition, Eigen is built on the standard software libraries for numerical linear algebra, BLAS and LAPACK, which are incredibly fast. These contribute greatly to the performance of the code. 

The code was developed in regards to three principal aims:
\begin{center}
	\begin{minipage}{0.2\textwidth}
		\begin{itemize}
			%\itemsep-0.6em
			\item fast,
			\item flexible,
			\item readable.
		\end{itemize}
	\end{minipage}
\end{center}
It needs to be flexible in order to support the Boltzmann machines as our trial wave function guess, and since we will try out various Jastrow factors it should be easy to add and remove wave function elements. Since quantum mechanical simulations in general are very expensive, it is important to develop efficient code to be able to study systems of some size. Lastly, we aim to write readable code such that others can reuse the code in its entirety or parts of it later. 

How we work to achieve the goals will be illustrated by code mainly picked from the \lstinline{WaveFunction} class, which is the heart of the code.

\section{Flexibility and legibility}
We have done several things in order to keep the code as legible as possible. Firstly, the code was written in an object orientated scheme which makes it more intuitive to a human as discussed in chapter \ref{chp:scientificprogramming}. Actually, the Hamiltonians, optimizers, wave functions, sampling methods and even the random number generator were treated as objects, making the code more or less as object orientated as possible. This makes it also easier to define a system, since we simply can set the preferred object. Below, we define a two-dimensional quantum dot system of 6 electrons with frequency $\omega=1.0$, learning rate $\eta=0.1$, number of Metropolis cycles $M=2^{20}=1048576$ and max number of iterations set to 1000.

\begin{lstlisting}[language={C++}]
System *QD = new System();

QD->setNumberOfDimensions(2);
QD->setNumberOfParticles(6);
QD->setNumberOfMetropolisSteps(int(pow(2, 20)));
QD->setFrequency(1.0);
QD->setLearningRate(0.1);

QD->setBasis(new Hermite(QD));
QD->setHamiltonian(new HarmonicOscillator(QD));

QD->setWaveFunctionElement(new Gaussian(QD));
QD->setWaveFunctionElement(new SlaterDeterminant(QD));
QD->setWaveFunctionElement(new PadeJastrow(QD));

QD->runIterations(1000);
\end{lstlisting}
We observe that the one first needs to define an object System, and then the various settings can be connected to this object. As you might notice, we use the \textbf{lowerCamelCase} naming convention for function and variable names, which means that each word begins with a capital letter except the initial word. For classes, we use the \textbf{UpperCamelCase} to distinguish from function names. This is known to be easy to read, and apart from for example the popular \textbf{snake\_case}, we do not need delimiters between the words, which saves some space. After the naming convention is decided, we are still responsible for giving reasonable names, which is not always an easy task, as Phil Karlton points out. When one sees the name, one should know exactly what the variable/function/class is or does. More about naming conventions can be read at reference \cite{noauthor_naming_2019}. In addition, for code format we use the ClangFormat, which provides a consequent way of formating the code. 

The snippet above also demonstrates how the code is made flexible when it comes to the wave function. One can simply construct a wave function consisting of various elements by calling the \lstinline{setWaveFunctionElement} multiple times. This creates a wave function vector, \lstinline{m_waveFunctionVector}, in the background containing all the elements. This makes it easy to compose the wave function in whatever way you want, all the elements can be combined. The reader might stubs on the use of the element \lstinline{Gaussian}, isn't the trial wave function defined by the Slater determinant multiplied with a Padé-Jastrow factor? It is, but as we will se later in section \ref{sec:factorizing}, the Gaussian part can be factorized out of the Slater determinant when using a Hermite basis. However, we will now start from the fundamental assumption that the trial wave function consists of a Slater determinant and a Jastrow factor, and take it from there. 

\section{Splitting the wave function in elements}
In a standard variational Monte Carlo computation, the trial wave function is assumed to consist of a single Slater determinant, and a Jastrow factor to take care of the repulsive interactions. Mathematically it can be expressed as
\begin{equation}
\Psi_T(\bs{r})=|\hat{D}(\bs{r})|J(\bs{r}_{ij})
\end{equation}
where the Slater determinant is the determinant of the matrix $\hat{D}(\bs{r})$, henceforth called the Slater matrix. 
To convince the reader that the Slater determinant and the Jastrow factor can be treated separately, we will consider a general trial wave function consisting of $p$ \textit{wave function elements} $\{\Phi_1, \Phi_2\hdots\Phi_p\}$,
\begin{empheq}[box={\mybluebox[5pt]}]{equation}
\Psi_T(\bs{r}) = \prod_{i=1}^p\Phi_i(\bs{r}),
\label{eq:elementproduct}
\end{empheq}
where $\bs{r}$ contains all the positions at a certain time. In order to treat each such element independently, we need to assume that a parameter $\theta_j$ only appears in an element.

\iffalse
In chapter \eqref{chp:quantum} we presented the basic principles behind a many-body trial wave function, including the Slater determinant and the well-known Padé-Jastrow factor. Further in chapter \eqref{chp:systems}, the common basis functions of the quantum dot and atomic systems were given, and in the previous chapter, \eqref{chp:machinelearning}, we explained how to create wave functions using Boltzmann machines. This means that all wave function elements used in this thesis already are presented, and in this chapter we collect them, together with their derivatives and various optimizations. The calculations below are based on two main assumptions:
\begin{enumerate}
	\item For each time step, we change one position coordinate only, i.e, move a particle along one of the principal axis.
	\item A variational parameter $\theta_i$ appears in only one of the wave function elements.
\end{enumerate}
The first assumption is useful when updating position dependent arrays. Typically, we only need to update an element in a vector or a row in a matrix when this assumption is raised, which is hugely beneficial with respect to the computational time. The last assumption makes all wave function elements independent, which obviously makes life easier. 

In this work, the trial wave function $\Psi_T$ is a product of all the $p$ wave function elements $\{\psi_1, \psi_2\hdots\psi_p\}$ that are involved in a calculation,
\begin{empheq}[box={\mybluebox[5pt]}]{equation}
\Psi_T(\bs{r}) = \prod_{i=1}^p\psi_i(\bs{r}),
\end{empheq}
where $\bs{r}$ contains all the positions at a certain time. 

As a practical example, consider a quantum dot of six interacting electrons. In that case, a Slater determinant needs to be included in order to ensure that the wave function is anti-symmetric, but since the Gaussian function appears in all the single particle function, this can be factorized out. We can therefore split the Slater determinant up in a simple Gaussian element and a determinant consisting of the Hermite polynomials. In addition, a Jastrow factor is required to take care of the correlations, so the trial wave function should consist of a total of three elements for this particular system:
\begin{equation*}
\Psi_T(\bs{r})=\psi_{sg}(\bs{r})\psi_{sd}(\bs{r})\psi_{jf}(\bs{r})
\end{equation*}
where $sg$ stands for the simple Gaussian, $sd$ stands for the Slater determinant and $jf$ stands for an arbitrary Jastrow factor. 

We will first try to convince the reader that the local energy and the parameter update can be calculated separately for each element, and then move on to find closed-form expressions for the required term in the calculations for all the elements.
\fi

\subsection{Kinetic energy computations} \label{sec:kinetic}
The local energy, defined in equation \eqref{eq:localenergy}, is
\begin{equation}
\begin{aligned}
E_L &=\frac{1}{\Psi_T(\bs{r})}\hat{\mathcal{H}}\Psi_T(\bs{r})\\
&=\sum_{k=1}^F\Big[-\frac{1}{2}\Big(\frac{1}{\Psi_T(\bs{r})}\nabla_k^2\Psi_T(\bs{r})\Big) + \mathcal{V}\Big].
\end{aligned}
\end{equation}
where we have $F=ND$ degrees of freedom. The first term, which is the kinetic energy term, is the only wave function-dependent one, and we will in this section split it up with respect to the elements. The potential energy term, $\mathcal{V}$, is not directly dependent on the wave function and will therefore not be further touched here. 

From the definition of differentiation of a logarithm, we have that
\begin{equation}
\frac{1}{\Psi_T(\bs{r})}\nabla_k\Psi_T(\bs{r})=\nabla_k\ln\Psi_T(\bs{r}),
\end{equation}
which provides the following useful relation 
\begin{equation}
\frac{1}{\Psi_T(\bs{r})}\nabla_k^2\Psi_T(\bs{r})=\nabla_k^2\ln\Psi_T(\bs{r}) + (\nabla_k\ln\Psi_T(\bs{r}))^2.
\end{equation}
Using the fact that the trial wave function is a product of all the elements, the term above is calculated by
\begin{equation}
\frac{1}{\Psi_T(\bs{r})}\nabla_k^2\Psi_T(\bs{r})=\sum_{i=1}^p\nabla_k^2\ln\Phi_i(\bs{r}) + \Big(\sum_{i=1}^p\nabla_k\ln\Phi_i(\bs{r})\Big)^2
\end{equation}
such that the total kinetic energy is given by
\begin{empheq}[box={\mybluebox[5pt]}]{equation}
-\frac{1}{2}\frac{1}{\Psi_T(\bs{r})}\nabla^2\Psi_T(\bs{r})=-\frac{1}{2}\bigg[\sum_{i=1}^p\nabla^2\ln\Phi_i(\bs{r}) + \sum_{k=1}^{F}\Big(\sum_{i=1}^p\nabla_k\ln\Phi_i(\bs{r})\Big)^2\bigg].
\label{eq:splittedkineticenergy}
\end{empheq}
This can be found when all local derivatives $\nabla^2\ln\Phi_i(\bs{r})$ and $\nabla_k\ln\Phi_i(\bs{r})$ are given. By assuming that the former is returned by a function \lstinline{computeLaplacian()} and the latter is returned by a function \lstinline{computeGradient(k)}, we compute the kinetic energy using the following function
\lstset{basicstyle=\scriptsize}
\begin{lstlisting}[language=c++,caption={From \lstinline{system.cpp}.}]
double System::getKineticEnergy()
{
	double kineticEnergy = 0;
	for (auto &i : m_waveFunctionElements) {
		kineticEnergy += i->computeLaplacian();
	}
	for (int k = 0; k < m_degreesOfFreedom; k++) {
		double nablaLnPsi = 0;
		for (auto &i : m_waveFunctionElements) {
			nablaLnPsi += i->computeGradient(k);
		}
		kineticEnergy += nablaLnPsi * nablaLnPsi;
	}
	return -0.5 * kineticEnergy;
}
\end{lstlisting}
Note that all the variables that are declared globally, here the vector \lstinline{m_waveFunctionElements} and the integer \lstinline{m_degreesOfFreedom} are denoted by a \lstinline{m_} to distinguish them from the variables declared locally. 

\subsection{Parameter gradients}
In section \ref{sec:parameterupdate}, we presented how the parameters can be updated by minimizing the energy expectation value. We recall that the only closed-form expression needed in addition to the local energy is $\nabla_{\theta_j}\ln\Psi_T(\bs{r}_i)$, which needs to be found. By applying equation \eqref{eq:elementproduct}, we find that
\begin{empheq}[box={\mybluebox[5pt]}]{equation}
\nabla_{\theta_j}\ln\Psi_T(\bs{r})=\sum_{i=1}^p\nabla_{\theta_j}\ln\Phi_i(\bs{r}),
\end{empheq}
which means that we need to find closed-form expressions of $\nabla_{\theta_j}\ln\Phi_i(\bs{r})$ for all wave function elements $\Phi_i(\bs{r})$ and all variational parameters $\theta_{j}$.

In the code, we store all the parameters in a \textit{parameter matrix} where each element has its own row of parameters. Similarly, we create a \textit{gradient matrix} of the same size to store the gradient $\nabla_{\theta_j}\ln\Phi_i(\bs{r})$ for each variational parameter. The implementation to get all the gradients are straight-forward, and given by
\begin{lstlisting}[language=c++,caption={From \lstinline{system.cpp}.}]
Eigen::MatrixXd System::getAllParameterGradients()
{
	for (int i = 0; i < m_numberOfElements; i++) {
		m_gradients.row(i) = m_waveFunctionElements[i]->computeParameterGradient();
	}
	return m_gradients;
}
\end{lstlisting}
where \lstinline{m_gradients} has the same number of rows as number of elements and the same number of columns as the maximum number of parameters in an element. The function \lstinline{computeParameterGradient()} returns a vector with all the gradients $\nabla_{\theta_j}\ln\Phi_i(\bs{r})$ of the respective element. Even though the gradients are used to update the parameters, we will postpone the discussion of the parameter update to section \ref{sec:update}.

\subsection{Probability ratio} \label{sec:probabilityratio}
How much a wave function element can be optimized heavily depends on the specific form of the element. For instance, sometimes the previous and present $\nabla_k\ln\Phi_i$ are closely related, and only differ from each other by a factor, while for some other elements they are not related at all. Those subjective optimizations will therefore be described when presenting each wave function element. 

However, there are still optimizations that apply to all elements and give great speed-up. An example is when calculating the ratio between the previous and present wave functions squared for all wave function elements instead of the wave function itself. Firstly, this is usually cheaper to calculate than the wave function itself because we are working in the logarithm space. Secondly, the ratio is actually what we use in the sampling, so it is a natural thing to calculate. The total probability ratio is just the product of all the probability ratios
\begin{empheq}[box={\mybluebox[5pt]}]{equation}
	\frac{P(\bs{r}_{\text{new}})}{P(\bs{r}_{\text{old}})}=\frac{|\Psi_T(\bs{r}_{\text{new}})|^2}{|\Psi_T(\bs{r}_{\text{old}})|^2}=\prod_{i=1}^p\frac{|\Phi_i(\bs{r}_{\text{new}})|^2}{|\Phi_i(\bs{r}_{\text{old}})|^2},
\end{empheq}
and below we will calculate this ratio for all the elements since we are going to use that directly in the sampling. We name the function returning the ratio for a particular element \lstinline{evaluateRatio()}, and we obtain the total probability ratio in the following way
\begin{lstlisting}[language=c++,caption={From \lstinline{system.cpp}.}]
double System::evaluateProbabilityRatio()
{
	double ratio = 1;
	for (auto &i : m_waveFunctionElements) {
		ratio *= i->evaluateRatio();
	}
	return ratio;
}
\end{lstlisting}

With this, we have introduced the four central functions of the wave function elements: \lstinline{computeLaplacian()}, \lstinline{computeGradient(k)}, \lstinline{computeParameterGradient()} and \lstinline{evaluateRatio()}. We will now proceed further and see how the trial wave function is split in various elements, and then find closed-form expressions to be used in these four functions for each element. Initially, the Slater determinant will be evaluated. 

\section{Slater determinant}
As we have seen above, the Slater determinant is the fundamental part of the trial wave function, in addition to the Jastrow factor. The main problem with the Slater determinant, is that it is very expensive to deal with as the number of particles increases. To find the gradient of the Slater determinant, as requested by equation \eqref{eq:splittedkineticenergy}, we need to compute the inverse of the Slater matrix, with by standard LU decomposition scales as $\sim\mathcal{O}(N^3)$ for a $N\times N$ matrix \cite{trahan_computational_2006}. Fortunately, there exist algorithms that lets us obtain the inverse of the matrix by recursive relations, scaling as $\sim\mathcal{O}(N^2)$. This will be detailed in section \ref{sec:efficientcalculationsofslaterdeterminant}.

Additionally, the Slater determinant can be split in a spin-up part and a spin-down part, which reduces the Slater matrix to two $N/2\times N/2$ matrices. Also factorizing common functions from the Slater determinant will give some speed-up. In this section, we will mostly discuss the various methods to make the update of the Slater matrix more efficient. 

The Slater determinant was first introduced in section \ref{sec:slater}, and looks like
\begin{equation}
\Phi_{sd}(\bs{r},\bs{\sigma})=|\hat{D}(\bs{r},\bs{\sigma})|\propto
\begin{vmatrix}
\psi_1(\boldsymbol{r}_1,\sigma_1) & \psi_2(\boldsymbol{r}_1,\sigma_2) & \hdots & \psi_N(\boldsymbol{r}_1,\sigma_N)\\
\psi_1(\boldsymbol{r}_2,\sigma_1) & \psi_2(\boldsymbol{r}_2,\sigma_2) & \hdots & \psi_N(\boldsymbol{r}_2,\sigma_N)\\
\vdots & \vdots & \ddots & \vdots \\
\psi_1(\boldsymbol{r}_N,\sigma_1) & \psi_2(\boldsymbol{r}_N,\sigma_2) & \hdots & \psi_N(\boldsymbol{r}_N,\sigma_N)
\end{vmatrix}
\label{eq:generalslater}
\end{equation}
where $\psi_j(\bs{r}_i,\sigma_j)$ is the single particle function occupying spot $i,j$ in the matrix.

\subsection{Splitting up the Slater determinant} \label{sec:splittingofslater}
In order to reduce the computational cost of the Slater determinant, we will split it up in a spin-up part and a spin-down part. In real life, we cannot immediately tell if an electron has spin-up or spin-down, but since we in the program need to know which particles that have spin-up and spin-down, we simply decide that the first $N_{\uparrow}$ particle have spin up and the remaining particles have spin down. 

In addition to reduce the computational cost, splitting up the Slater determinant also makes it possible to factorize out the spin-part from the spin-orbitals. Only the spatial part of the single particle functions contribute to the energies, such that we with advantage can get rid of the spin parts. Recall that the spin-orbitals can be written as 
\begin{equation}
\psi(\bs{r},\sigma)=\phi(\bs{r})\otimes\chi(\sigma) 
\end{equation}
where only the spatial part, $\phi(\bs{r})$, is dependent on the coordinates. As we in the ground state have double degeneracy, the spatial part will be the same for pairwise spin-up and spin-down particles, and we arrange them as
\begin{equation}
\psi_j(\bs{r}_i,\sigma_i)=
\begin{cases}
\phi_j(\bs{r}_i)\chi_{\uparrow}(\sigma_i)\quad\quad&\text{if } j<N_{\uparrow}\\
\phi_j(\bs{r}_i)\chi_{\downarrow}(\sigma_i)\quad\quad&\text{if } j\geq N_{\uparrow}.
\end{cases}
\end{equation}
Since the first $N_{\uparrow}$ particles have spin up, $\sigma_i=\uparrow\quad\forall\quad i\in\{1,2,...,N_{\uparrow}\}$ and the remaining have spin down, we can now set up the Slater determinant in equation \eqref{eq:generalslater} where each single particle function is split in a spatial part and a spin part,
\begin{equation}
\Phi_{sd}(\boldsymbol{r},\bs{\sigma})\propto
\begin{vmatrix}
\phi_1(\boldsymbol{r}_1)\chi_{\uparrow}(\uparrow) & \hdots & \phi_{N_{\uparrow}}(\boldsymbol{r}_1)\chi_{\uparrow}(\uparrow) & \phi_{1}(\boldsymbol{r}_1)\chi_{\downarrow}(\uparrow) & \hdots & \phi_{N_{\downarrow}}(\boldsymbol{r}_1)\chi_{\downarrow}(\uparrow)\\
\vdots & & \vdots & \vdots & & \vdots \\
\phi_1(\boldsymbol{r}_{N_{\uparrow}})\chi_{\uparrow}(\uparrow) & \hdots & \phi_{N_{\uparrow}}(\boldsymbol{r}_{N_{\uparrow}})\chi_{\uparrow}(\uparrow) & \phi_{1}(\boldsymbol{r}_{N_{\uparrow}})\chi_{\downarrow}(\uparrow) & \hdots & \phi_{N_{\downarrow}}(\boldsymbol{r}_{N_{\uparrow}})\chi_{\downarrow}(\uparrow)\\
\phi_1(\boldsymbol{r}_{N_{\uparrow}+1})\chi_{\uparrow}(\downarrow) & \hdots & \phi_{N_{\uparrow}}(\boldsymbol{r}_{N_{\uparrow}+1})\chi_{\uparrow}(\downarrow) & \phi_{1}(\boldsymbol{r}_{N_{\uparrow}+1})\chi_{\downarrow}(\downarrow) & \hdots & \phi_{N_{\downarrow}}(\boldsymbol{r}_{N_{\uparrow}+1})\chi_{\downarrow}(\downarrow)\\
\vdots & & \vdots & \vdots & & \vdots \\
\phi_1(\boldsymbol{r}_N)\chi_{\uparrow}(\downarrow) & \hdots & \phi_{N_{\uparrow}}(\boldsymbol{r}_N)\chi_{\uparrow}(\downarrow) & \phi_{1}(\boldsymbol{r}_N)\chi_{\downarrow}(\downarrow) & \hdots & \phi_{N_{\downarrow}}(\boldsymbol{r}_N)\chi_{\downarrow}(\downarrow)\\
\end{vmatrix}.
\end{equation}
We observe that the the spin-up particles sometimes occupy spin-down states, which they are not allowed to. Therefore, half of the elements become zero and the determinant can be further expressed as
\begin{equation}
\Phi_{sd}(\boldsymbol{r},\bs{\sigma})\propto
\begin{vmatrix}
\phi_1(\boldsymbol{r}_1)\chi_{\uparrow}(\uparrow) & \hdots & \phi_{N_{\uparrow}}(\boldsymbol{r}_1)\chi_{\uparrow}(\uparrow) & 0 & \hdots & 0\\
\vdots & & \vdots & \vdots & & \vdots \\
\phi_1(\boldsymbol{r}_{N_{\uparrow}})\chi_{\uparrow}(\uparrow) & \hdots & \phi_{N_{\uparrow}}(\boldsymbol{r}_{N_{\uparrow}})\chi_{\uparrow}(\uparrow) & 0 & \hdots & 0\\
0 & \hdots & 0 & \phi_{N_{\uparrow}+1}(\boldsymbol{r}_{N_{\uparrow}+1})\chi_{\downarrow}(\downarrow) & \hdots & \phi_{N}(\boldsymbol{r}_{N_{\uparrow}+1})\chi_{\downarrow}(\downarrow)\\
\vdots & & \vdots & \vdots & & \vdots \\
0 & \hdots & 0 & \phi_{N_{\uparrow}+1}(\boldsymbol{r}_N)\chi_{\downarrow}(\downarrow) & \hdots & \phi_{N}(\boldsymbol{r}_N)\chi_{\downarrow}(\downarrow)\\
\end{vmatrix}
\end{equation}
where the Slater matrix now is block diagonal! For a general block diagonal matrix, the determinant is given by the product of the determinant of each block
\begin{equation}
\Phi_{sd}(\boldsymbol{r},\bs{\sigma})\propto|\hat{D}_{\uparrow}(\bs{r}_{\uparrow},\sigma_{\uparrow})|\cdot |\hat{D}_{\downarrow}(\bs{r}_{\downarrow},\sigma_{\downarrow})|
\end{equation}
which can be seen by writing the total matrix as a product over all the block diagonal matrix. $\hat{D}_{\uparrow}$ is the matrix containing all spin-up states and $\hat{D}_{\downarrow}$ is the matrix containing all spin-down states. Since all elements in the respective matrices contain the same spin function, it can be factorized out,
\begin{equation}
\Phi_{sd}(\boldsymbol{r},\bs{\sigma})\propto|\hat{D}_{\uparrow}(\bs{r}_{\uparrow})|\{\chi_{\uparrow}(\sigma_1^{\uparrow})\chi_{\uparrow}(\sigma_2^{\uparrow})\hdots\chi_{\uparrow}(\sigma_{N_{\uparrow}}^{\uparrow})\}\cdot |\hat{D}_{\downarrow}(\bs{r}_{\downarrow})|\{\chi_{\downarrow}(\sigma_{N_{\uparrow+1}}^{\downarrow})\hdots\chi_{\downarrow}(\sigma_{N-1}^{\downarrow})\chi_{\downarrow}(\sigma_{N}^{\downarrow})\}
\end{equation}
and omitted in the future study since the energy is independent of spin. We are then left with the Slater determinant
\begin{empheq}[box={\mybluebox[5pt]}]{equation}
\Phi_{sd}(\boldsymbol{r})\propto|\hat{D}_{\uparrow}(\bs{r}_{\uparrow})|\cdot |\hat{D}_{\downarrow}(\bs{r}_{\downarrow})|
\end{empheq}
which is independent of spin, i.e, the matrices now consist of the spatial functions $\phi_j(\bs{r}_i)$ as the elements. 

It is also worth to notice that the size of the spin-up matrix is determined by the number of spin-up particles, and it is similar for the spin-down matrix. This means that the we can change the total spin $S$ by adjusting the relative sizes of the determinants. In the implementation, however, we stick to the magic quantum numbers, where $N_{\uparrow}=N_{\downarrow}$. 

This section was heavily inspired by Daniel Nissenbaum's doctoral dissertation, see appendix I in \cite{nissenbaum_stochastic_2008}.

\subsection{Factorizing out elements} \label{sec:factorizing}
We have now seen how the Slater determinant can be split up in a spin-up part and a spin-down part. Before we evaluate these determinants, we should try to make the elements of the Slater matrices as simple as possible to save computational time. If all the elements have the same factor, the computations will get much cheaper if the factor is factorized out of the matrix. How this is possible is easiest to see if we express the Slater determinant on a summation form,
\begin{equation}
\Phi_{sd}(\bs{r})\propto\sum_{p}(-1)^p\hat{P}\phi_1(\bs{r}_1)\phi_2(\bs{r}_2)\hdots\phi_N(\bs{r}_N),
\label{eq:slatersum}
\end{equation}
where the $\hat{P}$ is the permutation operator, permuting coordinates pairwise and the sum runs over all the possible permutations. If all the (spatial) single particle functions $\phi_j(\bs{r})_i$ can be split in two functions $f_j(\bs{r}_i)$ and $g(\bs{r}_i)$ where the latter is common for all the single particle functions,
\begin{equation}
\phi_j(\bs{r}_i)=f_j(\bs{r}_i)g(\bs{r}_i)
\end{equation}
the Slater determinant can be rewritten as
\begin{equation}
\begin{aligned}
\Phi_{sd}(\bs{r})&\propto\sum_{p}(-1)^p\hat{P}f_1(\bs{r}_1)g(\bs{r}_1)f_2(\bs{r}_2)g(\bs{r}_2)\hdots f_N(\bs{r}_N)g(\bs{r}_N)\\
&=g(\bs{r}_1)g(\bs{r}_2)\hdots g(\bs{r}_N)\sum_{p}(-1)^p\hat{P}f_1(\bs{r}_1)f_2(\bs{r}_2)\hdots f_N(\bs{r}_N)\\
&=\prod_{i=1}^Ng(\bs{r}_i)
\begin{vmatrix}
f_1(\boldsymbol{r}_1) & f_2(\boldsymbol{r}_1) & \hdots & f_N(\boldsymbol{r}_1)\\
f_1(\boldsymbol{r}_2) & f_2(\boldsymbol{r}_2) & \hdots & f_N(\boldsymbol{r}_2)\\
\vdots & \vdots & \ddots & \vdots \\
f_1(\boldsymbol{r}_N) & f_2(\boldsymbol{r}_N) & \hdots & f_N(\boldsymbol{r}_N)
\end{vmatrix}
\end{aligned}
\end{equation}
This is very useful when it comes to many common basises. For instance, the Hermite basis is given by 
\begin{equation}
\phi_j(\bs{r}_i)\propto H_j(\bs{r}_i)\exp(-\frac{1}{2}\omega|\bs{r}_i|^2)
\end{equation}
where $H_j(\bs{r}_i)$ are the Hermite polynomials and the Gaussian part fulfills the requirement of $g(\bs{r}_i)$. Therefore, we can construct a Slater determinant containing the Hermite polynomials only, treating the Gaussian as an independent element. This is not only preferable from an efficiency point of view, by doing this the variational parameter in the Gaussian is also removed from the determinant, which means that we can implement the determinant without worrying about the variational parameters. 

With this in mind, we will first treat the Gaussian element, obtaining its derivative and optimization schemes.  Moreover, in section \ref{sec:slaterdeterminant} we will discuss how the determinant can be treated efficiently. 

\subsection{Gaussian} \label{sec:simplegaussian}
When factorizing out the Gaussian part from the Slater determinant, we obtain the element
\begin{equation}
\Phi_{sg}(\bs{r}; \alpha)=\prod_{j=1}^Ng(\bs{r}_j)=\exp\Big(-\frac{1}{2}\omega\alpha\sum_{j=1}^Nr_j^2\Big)=\exp\Big(-\frac{1}{2}\omega\alpha\sum_{j=1}^{F}x_j^2\Big),
\end{equation}
with $N$ number of particles and $F=NP$ degrees of freedom. $\omega$ is the oscillator strength and $\alpha$ is a variational parameter, which for non-interacting atoms is 1. Because of the appearance of $r_i^2$, the function can easily be treated both in Cartesian and spherical coordinates, but in this thesis we will focus on the former.

We now use $j$ as our summation index, and reserve $i$ for the moved particle and $k$ as our the differentiating index. When changing a coordinate $x_i$ from $x_i^{\text{old}}$ to $x_i^{\text{new}}$, the probability ratio can easily be found to be 
\begin{empheq}[box={\mybluebox[5pt]}]{equation}
\label{eq:simplegaussianprobabilityratio}
\frac{|\Phi_{sg}(\bs{x}_{\text{new}})|^2}{|\Phi_{sg}(\bs{x}_{\text{old}})|^2}=\exp\Big(\omega\alpha\big((x_{i}^{\text{old}})^2-(x_{i}^{\text{new}})^2\big)\Big),
\end{empheq}
which contains only two elements! The gradient of $\ln\Phi_{\text{sg}}$ with respect to the coordinate $x_k$ is
\begin{empheq}[box={\mybluebox[5pt]}]{equation}
\nabla_k\ln\Phi_{sg}=-\omega\alpha x_k,
\end{empheq}
and the corresponding Laplacian is
\begin{empheq}[box={\mybluebox[5pt]}]{equation}
\nabla^2\ln\Phi_{sg}=-\omega\alpha F.
\end{empheq}
Finally, we will update $\alpha$ according to
\begin{empheq}[box={\mybluebox[5pt]}]{equation}
\nabla_{\alpha}\ln\Phi_{sg} = -\frac{1}{2}\omega\sum_{j=1}^Fx_j^2.
\end{empheq}
Since this wave function element is quite simple, there is no special optimization available that will cause a noticeable speed-up. The implementation looks like
\begin{lstlisting}[language=c++,caption={From \lstinline{gaussian.cpp}.}]
double Gaussian::evaluateRatio()
{
	return m_probabilityRatio;
}

double Gaussian::computeGradient(const int k)
{
	return -m_omega * m_alpha * m_positions(k);
}

double Gaussian::computeLaplacian()
{
	return -m_omega * m_alpha * m_degreesOfFreedom;
}

Eigen::VectorXd Gaussian::computeParameterGradient()
{
	Eigen::VectorXd gradients = Eigen::VectorXd::Zero(m_maxParameters);
	gradients(0) = -0.5 * m_omega * m_positions.cwiseAbs2().sum();
	return gradients;
}
\end{lstlisting}
where the probability ratio is calculated using
\begin{lstlisting}[language=c++,caption={From \lstinline{gaussian.cpp}.}]
double void Gaussian::updateProbabilityRatio(int changedCoord)
{
	m_probabilityRatio = exp(m_omega * m_alpha
				* (m_positionsOld(changedCoord) * m_positionsOld(changedCoord)
					- m_positions(changedCoord) * m_positions(changedCoord)));
}
\end{lstlisting}
We see that matrix-vector operations are used when it is possible, which makes the computations very efficient.

\iffalse
\subsection{Hydrogen-like orbitals} \label{sec:hydrogenlike}
The Hydrogen-like orbitals were presented in \eqref{eq:hydrogenlike}, but as we discussed earlier they cause some problems for atoms of the size of Neon and larger due to complex numbers. Instead, we decided to look at hydrogen-like orbitals with solid harmonics. Even though they do not have problems with complex numbers, they are quite complicated to differentiate, and the closed form will therefore be found by symbolic differentiating on the computer. However, we will do the exercise for the simplest case, which is sufficient for finding the Hydrogen and Helium ground states. This reads
\begin{equation}
\Phi_{hl}( \bs{r};\alpha)=\exp\Big(-Z\alpha\sum_{j=1}^Nr_j\Big)
\end{equation}
where $r_j$ is the distance from particle $j$ to the center. We then differentiate with respect to coordinate $x_k$, and obtain
\begin{equation}
\nabla_k\ln\Phi_{hl}=-Z\alpha\frac{x_k}{r_{k'}}
\end{equation}
The Laplacian is then given by
\begin{equation}
\nabla_k^2\ln\Phi_{hl}=-Z\alpha\Big(1-\frac{x_k^2}{r_{k'}^2}\Big)\frac{1}{r_{k'}}
\end{equation}
and the differentiation with respect to the variational parameter $\alpha$ is
\begin{equation}
\partial_{\alpha}\ln\Phi_{\text{hl}}=-Z\sum_{j=1}^Nr_j.
\end{equation}

For close-form expressions for higher order wave functions, please run the script \lstinline{generateHydrogenOrbitals.py}.
\fi

\subsection{The determinant} \label{sec:slaterdeterminant}
As discussed in section \ref{sec:splittingofslater}, the Slater determinant can be split in a spin-up part and a spin-down part, and further the common functions can be factorized out as shown in section \ref{sec:factorizing}. This means that the remaining determinant is not the full Slater determinant, and to distinguish it from the real Slater determinant, $\Phi_{sd}(\bs{r})$, we will name the element $\Phi_{det}(\bs{r})$. This determinant is course still split like Slater determinant, 
\begin{equation}
\Phi_{det}(\bs{r})=
|\hat{D}_{\uparrow}(\bs{r}_{\uparrow})|\cdot |\hat{D}_{\downarrow}(\bs{r}_{\downarrow})|,
\end{equation}
where $r_{\uparrow}$ are the coordinates of particles with spin up (defined as the first $N_{\uparrow}$ coordinates) and $r_{\downarrow}$ are the coordinates of particles with spin down (defined as the last $N_{\downarrow}$ coordinates). 

We can now utilize the logarithmic scale, by using that the logarithm of a product corresponds to summarize the logarithm of each factor,
\begin{equation}
\ln\Phi_{det}=\ln|\hat{D}_{\uparrow}(\bs{r}_{\uparrow})|+\ln|\hat{D}_{\downarrow}(\bs{r}_{\downarrow})|
\end{equation}
such that we only need to care about one of the determinants when differentiating, dependent on whether the coordinate we differentiate with respect to is among the spin-up or the spin-down coordinates:
\begin{equation}
\nabla_k\ln\Phi_{det}=
\begin{cases} 
\nabla_k\ln|\hat{D}_{\uparrow}(\bs{r}_{\uparrow})| & \text{if} \quad k<N_{\uparrow}\\
\nabla_k\ln|\hat{D}_{\downarrow}(\bs{r}_{\downarrow})| & \text{if} \quad k\geq N_{\uparrow}.
\end{cases}
\end{equation}
Before we go further, we will introduce a more general notation which cover both the cases:
\begin{equation}
\hat{D}(\bs{r})\equiv \hat{D}_{\sigma}(\bs{r}_{\sigma})
\end{equation}
where $\sigma$ is the spin projection. When summarizing, the sum is always over all relevant coordinates. Furthermore, we have that
\begin{equation}
\nabla_k\ln|\hat{D}(\bs{r})|=\frac{\nabla_k|\hat{D}(\bs{r})|}{|\hat{D}(\bs{r})|}
\end{equation}
and
\begin{equation}
\nabla_k^2\ln|\hat{D}(\bs{r})|=\frac{\nabla_k^2|\hat{D}(\bs{r})|}{|\hat{D}(\bs{r})|}-\bigg(\frac{\nabla_k|\hat{D}(\bs{r})|}{|\hat{D}(\bs{r})|}\bigg)^2
\end{equation}
At this point, there are (at least) two possible paths to the final expressions. We can keep on using matrix operations and find the expressions of $\nabla_k|\hat{D}(\bs{r})|/|\hat{D}(\bs{r})|$ and $\nabla_k^2|\hat{D}(\bs{r})|/|\hat{D}(\bs{r})|$ using Jacobi's formula, or we can switch to element representations of the matrices. We choose the latter, because we believe that gives the simplest derivations. 

The determinant of an arbitrary matrix $\hat{A}$ can be expressed by its comatrix $\hat{C}$ in the following way,
\begin{equation}
|\hat{A}|=\sum_{ij}a_{ij}c_{ji}
\end{equation}
where $a_{ij}$ are the matrix elements of $\hat{A}$ and $c_{ij}$ are the element of the comatrix. Going further, an element $c_{ij}$ can be expressed in terms of an element from the inverse of $\hat{A}$, $a_{ij}^{-1}$ \cite{morten_hjorth-jensen_computational_2019},
\begin{equation}
c_{ij}=a_{ij}^{-1}|\hat{A}|.
\end{equation}
Relating this to our particular problem, we can express 
\begin{equation}
\begin{aligned}
\frac{\nabla_k|\hat{D}|}{|\hat{D}|}&=\frac{\nabla_k\sum_{ij}d_{ij}c_{ji}}{\sum_{ij}d_{ij}c_{ji}}=\frac{\sum_j\nabla_kd_{kj}c_{jk}}{\sum_{ij}d_{ij}c_{ji}}\\
&=\frac{\sum_j\nabla_kd_{kj}d_{jk}^{-1}|\hat{D}|}{\sum_{ij}d_{ij}d_{ji}^{-1}|\hat{D}|}=\sum_j\nabla_kd_{kj}d_{jk}^{-1}
\end{aligned}
\label{eq:slaterelementshit}
\end{equation}
where $d_{ij}$ are the elements of $\hat{D}$ and we have used the fact that the elements $\nabla_kd_{ij}$ contribute to the sum if and only if $i=k$, such that the sum over $i$ collapses. Moreover, we use that multiplying a matrix with its inverse is identity, i.e, $\sum_{ij}d_{ij}d_{ji}^{-1}=1$ and the determinants cancel. Similarly, we get 
\begin{equation}
\frac{\nabla_k^2|\hat{D}|}{|\hat{D}|}=\sum_j\nabla_k^2d_{kj}d_{jk}^{-1}
\end{equation}
for the Laplacian. We are then set to write up the final expressions for the gradient and Laplacian of the logarithm of the Slater determinant,
\begin{equation}
\begin{aligned}
\nabla_k\ln|\hat{D}(\bs{r})|&=\sum_{j}d_{jk}^{-1}(\bs{r})\nabla_k\phi_{j}(\bs{r}_k)\\
\nabla_k^2\ln|\hat{D}(\bs{r})|&=\sum_jd_{jk}^{-1}(\bs{r})\nabla_k^2\phi_{j}(\bs{r}_k)-\Big(\sum_jd_{jk}^{-1}(\bs{r})\nabla_k\phi_{j}(\bs{r}_k)\Big)^2
\end{aligned}
\label{eq:derivativelndet}
\end{equation}
where we have used that $d_{ij}=\phi_j(\bs{r}_i)$ with $\phi_j(\bs{r}_i)$ as a single particle function found in the Slater determinant, see above.

\subsubsection{Efficient calculation of Slater determinants} \label{sec:efficientcalculationsofslaterdeterminant}
As aforementioned, dealing with the Slater determinant is very computational expensive, mainly because of the requirement of the inverse Slater matrix. However, by exploiting that only one row in the Slater matrix is updated for each step, we can update the inverse recursively. By using the same element representations as above, the ratio between the new and the old determinant can be expressed as
\begin{equation}
R\equiv \frac{|\hat{D}(\bs{r}_{\text{new}})|}{|\hat{D}(\bs{r}_{\text{old}})|}=\frac{\sum_{j}d_{ij}(\bs{r}_{\text{new}})c_{ij}(\bs{r}_{\text{new}})}{\sum_{j}d_{ij}(\bs{r}_{\text{old}})c_{ij}(\bs{r}_{\text{old}})}=\sum_{j}d_{ij}(\bs{r}_{\text{new}})d_{ji}^{-1}(\bs{r}_{\text{old}})
\label{eq:slaterratio}
\end{equation}
which is very similar to the calculation given in equation \eqref{eq:slaterelementshit}. To calculate the inverse matrix $\hat{D}^{-1}$ efficiently, we need to calculate
\begin{equation}
S_j=\sum_{l=1}^Nd_{il}(\bs{r}_{\text{new}})d_{lj}^{-1}(\bs{r}_{\text{old}})
\label{eq:slaters}
\end{equation}
for all columns but the one associated with the moved particle, $i$. For all columns where $j\neq i$, we then find the new elements using
\begin{equation}
d_{kj}^{-1}(\bs{r}_{\text{new}})=d_{kj}^{-1}(\bs{r}_{\text{old}})-\frac{S_j}{R}d_{ki}^{-1}(\bs{r}_{\text{old}})
\label{eq:slaterinverse}
\end{equation}
while the remaining column, $i$, is updated using the simple formula \cite{morten_hjorth-jensen_computational_2019}
\begin{equation}
d_{ki}^{-1}(\bs{r}_{\text{new}})=\frac{1}{R}d_{ki}^{-1}(\bs{r}_{\text{old}}).
\label{eq:slaterinverse2}
\end{equation}
Those procedures makes the inverting scale as $\mathcal{O}(N^2)$ instead of $\mathcal{O}(N^3)$, which is largely beneficial for large systems.

We assume that we do not have any variational parameter in the Slater determinant, and obtain three expressions of the case when a particle with spin up is moved and three of the case when a particle with spin down is moved. 

\begin{empheq}[box={\mybluebox[5pt]}]{equation}
\begin{aligned}
&\quad\text{if}\quad k<N_{\uparrow}:\\
\frac{|\Phi_{sd}(\bs{r}_{\text{new}})|^2}{|\Phi_{sd}(\bs{r}_{\text{old}})|^2}&=
\frac{|\hat{D}_{\uparrow}(\bs{r}_{\uparrow}^{\text{new}})|^2}{|\hat{D}_{\uparrow}(\bs{r}_{\uparrow}^{\text{old}})|^2}\\
\nabla_k\ln|\hat{D}_{\uparrow}(\bs{r}_{\uparrow})|&=\sum_{j=1}^{N_{\uparrow}}\nabla_kd_{jk}(\bs{r}_{\uparrow})d_{kj}^{-1}(\bs{r}_{\uparrow})
\end{aligned}
\label{eq:slaterupdateup}
\end{empheq}

\begin{empheq}[box={\mybluebox[5pt]}]{equation}
\begin{aligned}
&\quad\text{if}\quad k\geq N_{\uparrow}:\\
\frac{|\Phi_{sd}(\bs{r}_{\text{new}})|^2}{|\Phi_{sd}(\bs{r}_{\text{old}})|^2}&=
\frac{|\hat{D}_{\downarrow}(\bs{r}_{\downarrow}^{\text{new}})|^2}{|\hat{D}_{\downarrow}(\bs{r}_{\downarrow}^{\text{old}})|^2}\\
\nabla_k\ln|\hat{D}_{\downarrow}(\bs{r}_{\downarrow})|&=\sum_{j=N_{\uparrow}}^{F}\nabla_kd_{jk}(\bs{r}_{\downarrow})d_{kj}^{-1}(\bs{r}_{\downarrow})
\end{aligned}
\end{empheq}

\begin{empheq}[box={\mybluebox[5pt]}]{equation}
\nabla^2\ln|\hat{D}(\bs{r})|=\sum_{k=1}^F\bigg[\sum_{j=1}^{F}\nabla_k^2d_{jk}(\bs{r})d_{kj}^{-1}(\bs{r})-\Big(\sum_{j=1}^{F}\nabla_kd_{ik}(\bs{r})d_{ki}^{-1}(\bs{r})\Big)^2\bigg]
\label{eq:slaterlaplacian}
\end{empheq}

We have now presented the theory behind finding the ratio between the new and the old probability and the gradients of the determinant (equation (\ref{eq:slaterupdateup}-\ref{eq:slaterlaplacian})), and we have described how we efficiently can find the inverse of the Slater matrix (equation(\ref{eq:slaterratio}-\ref{eq:slaterinverse2})). However, to make things more clear we will outline some selected parts of the Slater determinant implementation.

In the code, we create a Slater matrix, \lstinline{m_slaterMatrix} where all the elements are stored. This matrix contains both $\hat{D}_{\uparrow}$ and $\hat{D}_{\downarrow}$, and has therefore dimensions $N\times N/2$. Furthermore, we store the gradient of the elements with respect to all the $F$ elements in a matrix \lstinline{m_slaterMatrixDer}, which naturally gets the dimensions $F\times N/2$. We also create a matrix \lstinline{m_slaterMatrixSecDer} to store all the Laplacians of the elements which also has the dimensions $F\times N/2$. In all of them, we only need to update a row when a particle is moved, which makes it quite efficient. The \lstinline{m_slaterMatrixDer} is updated in the following way

\begin{lstlisting}[language={c++},caption={Taken from \lstinline{slaterdeterminant.cpp}.}]
void SlaterDeterminant::updateSlaterMatrixDerRow(const int row)
{
	int particle = int(row / m_numberOfDimensions);
	int dimension = row % m_numberOfDimensions;
	for (int col = 0; col < m_numberOfParticlesHalf; col++) {
		m_slaterMatrixDer(row, col) = m_basis->basisElementDer(col,
							dimension, m_positions.col(particle));
	}
}
\end{lstlisting}
where each element is taken from the \lstinline{basisElementDer} function in the \lstinline{Basis} class. This function returns just the derivative of the single particle function called for the chosen basis set. Note also that only the coordinates of the moved particle, stored in a column of the \lstinline{m_positions} matrix, is needed for the update. The update of \lstinline{m_slaterMatrix} and \lstinline{m_slaterMatrixSecDer} are very straight-forward and similar to the example above, so we will not discuss them further.

Something that might be less intuitive, is how to update the inverse of the Slater matrix. We also store this in a dedicated matrix \lstinline{m_slaterMatrixInverse}, and we use LU decomposition only to initialize it. Thereafter, we use the formulas above to update the inverse. We implement it as

\begin{lstlisting}[language={c++},caption={Taken from \lstinline{slaterdeterminant.cpp}.}]
void SlaterDeterminant::updateRatio()
{
	m_ratio = m_slaterMatrix.row(m_particle) * m_slaterMatrixInverse.col(m_particle);
}

void SlaterDeterminant::updateSlaterMatrixInverse(int start, int end)
{
	updateRatio();
	for (int j = start; j < m_particle; j++) {
		double S = m_slaterMatrix.row(m_particle) * m_slaterMatrixInverse.col(j);
		m_slaterMatrixInverse.col(j) -= S * m_slaterMatrixInverse.col(m_particle) 
											/ m_ratio;
	}
	for (int j = m_particle+1; j < end; j++) {
		double S = m_slaterMatrix.row(m_particle) * m_slaterMatrixInverse.col(j);
		m_slaterMatrixInverse.col(j) -= S * m_slaterMatrixInverse.col(m_particle) 
											/ m_ratio;
	}
	m_slaterMatrixInverse.col(m_particle) /= m_ratio;
}
\end{lstlisting}
where \lstinline{m_ratio} is a global variable also returned by the function \lstinline{evaluateRatio} (see section \ref{sec:optimization}). Note that the loops never affect the $i$'th columns, where particle $i$ is moved (in the code denoted by the global variable \lstinline{m_particle}). The arguments to the function \lstinline{updateSlaterMatrixInverse} specify which part of the matrix that should be updated, based on whether the moved particle has spin-up or spin-down.

We will end our discussions of the Slater determinant by presenting the implementation of the gradient and the Laplacian of the logarithm of the determinant. These were decided to be stored in the vectors \lstinline{m_determinantDer} and \lstinline{m_determinantSecDer} for $(\nabla_k|\hat{D}(\bs{r})|)/|\hat{D}(\bs{r})|$ and $(\nabla_k^2|\hat{D}(\bs{r})|)/|\hat{D}(\bs{r})|$ respectively. These vectors are updated using vector operations in the following fashion
\begin{lstlisting}[language={c++},caption={Taken from \lstinline{slaterdeterminant.cpp}.}]
void SlaterDeterminant::updateSlaterDeterminantDerivatives(int start, int end)
{
	for (int i = start * m_numberOfDimensions; i < end * m_numberOfDimensions; i++) {
		int particle = int(i / m_numberOfDimensions);
		m_determinantDer(i) = m_slaterMatrixDer.row(i) 
		    * m_slaterMatrixInverse.col(particle);
		m_determinantSecDer(i) = m_slaterMatrixSecDer.row(i)
			* m_slaterMatrixInverse.col(particle);
	}
}
\end{lstlisting}
We avoid a double loop by taking a inner product instead. However, we are left with one loop which can also be avoided using smart matrix operations.

\section{Jastrow factor}
The second main part of the trial wave function is the Jastrow factor, which is meant to take care of the electron-electron correlations. The optimization scheme of this element is not as complex as the Slater determinant, and this section will therefore be notably shorter than the previous. We will first discuss the two Jastrow factors given in section \ref{sec:jastrow}: the simple Jastrow and the Padé-Jastrow factor, and then we look at how the distance matrix can be updated efficiently. 

\subsection{Simple Jastrow factor}
Recall the simple Jastrow factor from \eqref{eq:SimpleJastrow},
\begin{equation}
\Phi_{sj}(\bs{r};\bs{\beta})=\exp\bigg(\sum_{i=1}^N\sum_{j>i}^N\beta_{ij}r_{ij}\bigg).
\end{equation}
with $N$ as the number of particles, $r_{ij}$ as the distance between particle $i$ and $j$ and $\beta_{ij}$ as variational parameters.

This is a quite simple element, but one challenge is that we operate in Cartesian coordinates, while the expressed Jastrow factor obviously is easier to deal with in polar coordinates. Since we need to differentiate this with respect to all degrees of freedom, we need to be attentive not confusing the particle indices with the coordinate indices. Let us reserve $j$ as the coordinate index and $j'$ as the index of the corresponding particle. The relationship between $j$ and $j'$ is \textit{always} $j'=j\setminus D$, where the backslash denotes integer division. The other way around, we have $j=j'+d$ where $d$ is the respective dimension of the coordinate $j$. With that notation, the probability ratio is given by
\begin{empheq}[box={\mybluebox[5pt]}]{equation}
\frac{|\Phi_{sj}(\bs{r}_{\text{new}})|^2}{|\Phi_{\text{sj}}(\bs{r}_{\text{old}})|^2}=\exp\bigg(2\sum_{j'=1}^N\beta_{i'j'}(r_{i'j'}^{\text{new}}-r_{i'j'}^{\text{old}})\bigg)
\end{empheq}
where $i'$ again is the moved particle. The gradient is straight-forward to find, and reads
\begin{empheq}[box={\mybluebox[5pt]}]{equation}
\nabla_k\ln\Phi_{sj}=\sum_{j'=1}^N\beta_{k'j'}\frac{x_k-x_j}{r_{k'j'}}
\end{empheq}
where $j$ is related to the same dimension as $k$. Here we use $x_j$ as a general coordinate, no matter whether it is associated with the x-direction or not. This also applies for the Laplacian,
\begin{empheq}[box={\mybluebox[5pt]}]{equation}
\nabla^2\ln\Phi_{sj}=\sum_{k=1}^{F}\sum_{j'=1}^N\frac{\beta_{k'j'}}{r_{k'j'}}\bigg[1-\Big(\frac{x_k-x_j}{r_{k'j'}}\Big)^2\bigg].
\end{empheq}
Finally, the parameter update is given by
\begin{empheq}[box={\mybluebox[5pt]}]{equation}
\nabla_{\beta_{m'l'}}\ln\Phi_{sj}=r_{m'l'}.
\end{empheq}
For this element, the most important thing we can do to keep the computational cost as low as possible is to reveal that only a row and a column of the distance matrix is changed as we change a coordinate. Updating the entire distance matrix means updating $N^2$ elements, while updating a row and a column means updating $2N$ elements, which is an essential difference for large systems. This is detailed in section \ref{sec:distancematrix}.

We also observe that the factor $(x_k-x_j)/r_{k'j'}$ is found in both the gradient and the Laplacian, so by storing this matrix we can speed-up the computations. Most naturally, the matrix has dimensions $F\times F$, but using that we only are interested in the elements where $x_k$ and $x_j$ have the same dimension and that the diagonal is zero, only $F$ of the elements need to be found. Further, the matrix is obviously anti-symmetric, so we actually only need to calculate $F/2$ of the elements. 

When one particle is moved, only the elements related to the moved particle need to be updated, which is $2(N-1)$ elements. Again, we utilize that the matrix is anti-symmetric and get the following efficient update scheme
\begin{lstlisting}[language={c++},caption={Taken from \lstinline{simplejastrow.cpp}.}]
void SimpleJastrow::updatePrincipalDistance(int i)
{
    int i_d = i % m_numberOfDimensions;
	for (int j_p = 0; j_p < i_p; j_p++) {
		int j = i_d + j_p * m_numberOfDimensions;
		m_principalDistance(i, j) = (m_positions(i) - m_positions(j)) 
										/ m_distanceMatrix(i_p, j_p);
		m_principalDistance(j, i) = -m_principalDistance(i, j);
	}
	for (int j_p = i_p + 1; j_p < m_numberOfParticles; j_p++) {
		int j = i_d + j_p * m_numberOfDimensions;
		m_principalDistance(i, j) = (m_positions(i) - m_positions(j)) 
										/ m_distanceMatrix(i_p, j_p);
		m_principalDistance(j, i) = -m_principalDistance(i, j);
	}
}
\end{lstlisting}
where \lstinline{i_p} is the moved particle \lstinline{i_d} is the dimension the particle is moved in and \lstinline|i| thus is the coordinate index. Similarly, the loop goes over the particles \lstinline{j_p} with the associated coordinate \lstinline{j}. Note that we split the loop in two parts to avoid calculating the distance from a particle to itself. This trick is also done in many of the other functions in the simple Jastrow class, and also in the Padé-Jastrow factor class. 

\subsection{The Padé-Jastrow factor}
The Padé-Jastrow factor is a more complicated Jastrow factor, and was specified in equation \eqref{eq:PadeJastrow}, 
\begin{equation}
\Phi_{pj}(\bs{r}; \beta) = \exp\bigg(\sum_{i=1}^N\sum_{j>i}^N\frac{a_{ij}r_{ij}}{1+\beta r_{ij}}\bigg).
\end{equation}
where $\beta$ is a variational parameter and the $a_{ij}$ is \textbf{not} a variational parameter, but rather constants dependent on the spin of particles $i$ and $j$.

Similarly to the simple Jastrow, we also here need to distinguish between particle indices and coordinate indices because of the radial distances $r_{ij}$. We do the same trick as presented for the simple Jastrow denoting $j$ as the coordinate index and $j'$ as the particle index, and obtain the gradient 
\begin{equation}
\nabla_k\ln\Phi_{pj}=\sum_{j'\neq k'=1}^N\frac{a_{k'j'}}{(1+\beta r_{k'j'})^2}\frac{x_k-x_j}{r_{k'j'}}
\end{equation}
with respect to the coordinate $x_k$. By again differentiating this with respect to $x_k$, we obtain the Laplacian
\begin{equation}
\nabla^2\ln\Phi_{pj}=\sum_{k=1}^{F}\sum_{j'\neq k'=1}^N\frac{a_{k'j'}}{(1+\beta r_{k'j'})^2}\bigg[1-\Big(1+2\frac{\beta r_{k'j'}}{1+\beta r_{k'j'}}\Big)\Big(\frac{x_k-x_j}{r_{k'j'}}\Big)^2\bigg]\frac{1}{r_{k'j'}}.
\end{equation}
Similar to the simple Jastrow factor, we again observe the factor $(x_k-x_j)/r_{k'j'}$ in both the gradient and the Laplacian, which can be stored as a matrix and updated in the same way as described for the simple Jastrow factor. The last expression we need  is the one used to update the variational parameter $\beta$, which is found to be
\begin{equation}
\nabla_{\beta}\ln\Phi_{pj}=-\sum_{i'=1}^N\sum_{j'>i'}^N\frac{a_{ij}r_{ij}^2}{(1+\beta r_{ij})^2}.
\end{equation}

In addition to the factor $g_{ij}\equiv(x_k-x_j)/r_{k'j'}$, there are multiple factors that we can store to make the computations cheaper. The factor $f_{ij}\equiv a_{ij}/(1+\beta r_{ij})^2$ is found both in the gradient, Laplacian and parameter gradient, and storing it will saving some significant computations. Lastly, the factor $h_{ij}\equiv r_{ij}/(1+\beta r_{ij})$ is found in several places and will be stored as well. As a summary, we use
\begin{equation}
f_{ij}=\frac{a_{ij}}{(1+\beta r_{ij})^2}\quad\quad g_{ij}=\frac{x_i-x_j}{r_{i'j'}}\quad\quad h_{ij}=\frac{r_{ij}}{1+\beta r_{ij}}.
\end{equation}
and obtain the simplified expressions
\begin{empheq}[box={\mybluebox[5pt]}]{equation}
\begin{aligned}
\frac{|\Phi_{pj}(\bs{r}_{\text{new}})|^2}{|\Phi_{pj}(\bs{r}_{\text{old}})|^2}&=\exp\Big(2\sum_{j'=1}^Na_{i'j'}(h_{i'j'}^{\text{new}}-h_{i'j'}^{\text{old}})\Big)\\
\nabla_k\ln\Phi_{pj} &=\sum_{j'\neq k'=1}^Nf_{k'j'}\cdot g_{kj}\\
\nabla^2\ln\Phi_{pj} &= \sum_{k=1}^F\sum_{j'\neq k'=1}^N\frac{f_{k'j'}}{r_{k'j'}}\Big[1-(1+2\beta h_{k'j'})g_{kj}^2\Big]\\
\nabla_{\beta}\ln\Phi_{pj}&=-\sum_{l'=1}^N\sum_{j>l}^Na_{l'j'}h_{l'j'}^2=-\sum_{l'=1}^N\sum_{j>l}^Nf_{l'j'}r_{l'j'}^2
\end{aligned}
\end{empheq}
with marked indices ($j'$) as the particle related ones and the unmarked ($j$) as the coordinate related ones. $i'$ is the moved particle. We now proceed further to the update of the distance matrix, which is where we can find the remaining optimization possibilities. 

\subsection{Updating the distance matrix} \label{sec:distancematrix}
The distance matrix, which is used in the Jastrow factors, gives an illustrating example on how we can avoid repeating calculations. The matrix, henceforth named $M$, contains the relative distances between all the particles, for three particles given by
\begin{eqnarray}
M=
\begin{pmatrix}
r_{11} & r_{12} & r_{13} \\
r_{21} & r_{22} & r_{23} \\
r_{31} & r_{32} & r_{33}
\end{pmatrix}
=
\begin{pmatrix}
0 & r_{12} & r_{13} \\
r_{12} & 0 & r_{23} \\
r_{13} & r_{23} & 0
\end{pmatrix}
\end{eqnarray}
where $r_{ij}$ means the distance between particles $i$ and $j$. Since $r_{ij}=r_{ji}$ and $r_{ii}=0$, the matrix becomes symmetric with zeros on the diagonal, which means that we only need to calculate $N(N-1)/2$ elements instead of $N^2$. Further, we can utilize that only a particle is moved at a time, which means that only a row and a column are changed when a particle is moved. For instance, if particle 1 is moved, the upper row and the left-hand-side column in matrix $M$ need to be updated. In our program, we have implemented this in the following way
\lstset{basicstyle=\scriptsize}
\begin{lstlisting}[language=c++,caption={Taken from \lstinline{metropolis.cpp}.}]
double Metropolis::calculateDistanceMatrixElement(const int i, const int j) 
{
	double dist = 0;
	int parti   = m_numberOfDimensions*i;
	int partj   = m_numberOfDimensions*j;
	for(int d=0; d<m_numberOfDimensions; d++) {
		double diff = m_positions(parti+d)-m_positions(partj+d);
		dist += diff*diff;
	}
	return sqrt(dist);
}

void Metropolis::calculateDistanceMatrixCross(const int particle) {
	for(int i=0; i<m_numberOfParticles; i++) {
		m_distanceMatrix(particle, i) = calculateDistanceMatrixElement(particle, i);
		m_distanceMatrix(i, particle) = m_distanceMatrix(particle, i);
	}
}
\end{lstlisting}
where the function \lstinline{calculateDistanceMatrixElement(i,j)} returns element \lstinline{i,j} of the matrix, which is called from the function \lstinline{calculateDistanceMatrixCross(particle)}. The latter takes the moved particle index as input, and updates the necessary row and column of the matrix. 

For systems of non-interacting particles, the distance matrix is redundant, and should therefore not be calculated. We have solved this by giving all the wave function elements and the Hamiltonians a number which indicated whether they require the distance matrix or not, as mentioned above. If no part of the code needs the distance matrix, it is never calculated. 

We also calculate the radial position globally when it is required by any part of the code. The components are stored in a vector named \lstinline{radialVector}, applying the same optimization ideas as the distance matrix. 

\section{Sampling} \label{sec:sampling}
Also when it comes to the sampling itself, there exist optimization schemes to speed-up the process. Remember that the sampling algorithm often is repeated millions of times for each iteration, so even a small impact can give a massive speed-up. We will initially present the brute force sampling implementation in its entirety, before we move on to the importance sampling implementation. For the latter, we will discuss the optimization possibilities and connect them to the actual implementation.

\subsection{Brute force sampling}
The brute force sampling was introduced in section \ref{sec:bruteforce}, and is the most basic sampling method implemented. The sampling function \lstinline{BruteForce::acceptMove}, which returns true if the move is accepted, is implemented as

\begin{lstlisting}[language=c++,caption={Taken from \lstinline{bruteforce.cpp}.}]
bool BruteForce::acceptMove()
{
	int i = m_RNG->nextInt(m_degreesOfFreedom);

	m_positionsOld = m_positions;
	m_radialVectorOld = m_radialVector;
	m_distanceMatrixOld = m_distanceMatrix;

	m_positions(i) += (m_RNG->nextDouble() - 0.5) * m_stepLength;
	if (m_calculateDistanceMatrix) {
		Metropolis::calculateDistanceMatrixCross(int(i / m_numberOfDimensions));
	}
	if (m_calculateRadialVector) {
		Metropolis::calculateRadialVectorElement(int(i / m_numberOfDimensions));
	}
	m_system->updateAllArrays(m_positions, m_radialVector, m_distanceMatrix, i);

	double p = m_system->evaluateProbabilityRatio();
	if (p < m_RNG->nextDouble()) {
		m_positions = m_positionsOld;
		m_distanceMatrix = m_distanceMatrixOld;
		m_radialVector = m_radialVectorOld;
		m_system->resetAllArrays();
		return false;
	}
return true;
}
\end{lstlisting}
where \lstinline|i| is the changed coordinate which is drawn from the random number generator \lstinline|m_RNG|. Initially the old positions, radial vector and distance matrix are stored in case the move is rejected, and then a new move is proposed in positive or negative direction. If the radial vector or distance matrix (or both) are needed somewhere in the code, they are updated in this function, using the ideas and implementation presented in section \ref{sec:distancematrix}. The they are distributed to the wave function elements using the function \lstinline|updateAllArrays|.

In the end, the probability ratio is evaluated using the function \lstinline|evaluateProbabilityRatio| presented in section \ref{sec:probabilityratio}. If this ratio is larger than a random number between 0 and 1, the move is accepted, and otherwise we set all the arrays back to the old ones (also the ones in the wave function elements). 

\subsection{Importance sampling}
The importance sampling implementation is very similar to the brute force sampling implementation, and we will therefore not repeat it. However, we need to calculate the quantum force and the ratio between the new and the old Green's function, which can be calculated in clever ways to keep the code efficient. 

We have already sen that the quantum force takes the same form as the gradient of the trial wave function, $\bs{F}(\bs{r})=2(\nabla\Psi_T(\bs{r}))/\Psi_T(\bs{r})$, and we can therefore simply reuse the function \lstinline|computeGradient|, which is a part of the local energy computations from section \ref{sec:kinetic}. We call this from the function \lstinline|ImpotanceSampling::QuantumForce|, which contains the few lines of code

\begin{lstlisting}[language=c++,caption={Taken from \lstinline{importancesampling.cpp}.}]
double ImportanceSampling::QuantumForce(const int i)
{
	double QF = 0;
	for (auto &j : m_waveFunctionVector) {
		QF += j->computeGradient(i);
	}
	return 2 * QF;
}
\end{lstlisting}
where the force in dimension \lstinline|i| is returned. 

The Green's function was first presented in section \ref{sec:importancesampling}, and at first glance it might look computational expensive to evaluate. Fortunately, we only need the ratio between the old and the new function which can be found in a quite simple fashion. Actually, both the diffusion constant $D$ and the time step $\Delta t$ cancel in the exponent, and the ratio can be expressed in the elegant form
\begin{equation}
g(\bs{r}',\bs{r},\Delta t)\equiv\frac{G(\bs{r}',\bs{r},\Delta t)}{G(\bs{r},\bs{r}',\Delta t)}=\exp((\bs{r}'-\bs{r})\cdot(\bs{F}(\bs{r})-\bs{F}(\bs{r}'))/2)
\label{eq:greensratio}
\end{equation}
where $\bs{r}$ and $\bs{r}'$ differ by one element and so does $\bs{F}(\bs{r})$ and $\bs{F}(\bs{r}')$. It can therefore be evaluated 

\iffalse
\IncMargin{1em}
\begin{algorithm}
	\SetAlgoLined
	\Parameter{$\Delta t$: Fictive time step}
	\Data{$\bs{r}'$: Initial particle positions}
	\Require{$\Psi_T(\bs{r})$: Initial trial wave function guess}
	\BlankLine
	$\bs{F}(\bs{r}')\leftarrow2(\nabla\Psi_T(\bs{r}'))/\Psi_T(\bs{r}')$ (Initialize the quantum force) \;
	$p(\bs{r}',\bs{r})\leftarrow1$ (Initialize the probability ratio) \;
	$g(\bs{r}',\bs{r},\Delta t)\leftarrow1$ (Initialize Green's function ratio) \;
	\For{$i\leftarrow 1$ \KwTo $M$}{
		$\bs{r}\leftarrow\bs{r}'$ (Save current positions $\bs{r}'$ in a vector $\bs{r}$) \;
		$\bs{F}(\bs{r})\leftarrow\bs{F}(\bs{r}')$ (Save the current quantum force, $\bs{F}(\bs{r}')$, in a vector $\bs{F}(\bs{r})$) \;
		$p(\bs{r},\bs{r}')\leftarrow p(\bs{r}',\bs{r})$ (Save current probability ratio)\;
		$g(\bs{r},\bs{r}',\Delta t)\leftarrow g(\bs{r}',\bs{r},\Delta t)$ (Save the current Green's function ratio) \;
		\BlankLine
		$\bs{r}'\leftarrow\bs{r}+D\Delta t\bs{F}(\bs{r}) + \bs{\chi}\sqrt{\Delta t}$ (Update position based on the Langevin equation) \;
		$\bs{F}(\bs{r}')\leftarrow2\nabla\ln\Psi_T(\bs{r'})$ (Update the quantum force) \;
		$p(\bs{r}',\bs{r})=|\Psi_T(\bs{r}')|^2/|\Psi_T(\bs{r})|^2$ (Update the probability ratio) \;
		$g(\bs{r}',\bs{r},\Delta t)=G(\bs{r}',\bs{r},\Delta t)/G(\bs{r},\bs{r}',\Delta t)$ (Update the Green's function ratio) \;
		$w\leftarrow p(\bs{r}',\bs{r}) g(\bs{r}',\bs{r},\Delta t)$ (Calculate acceptance probability) \;
		$q\leftarrow\mathcal{U}(0,1)$ (Draw a random number between 0 and 1)\;
		\eIf{$w<q$}{
			$\bs{r}'\leftarrow\bs{r}$ (Reset positions)\;
			$\bs{F}(\bs{r}')\leftarrow\bs{F}(\bs{r})$ (Reset the quantum force) \;
			$p(\bs{r}',\bs{r})\leftarrow p(\bs{r},\bs{r}')$ (Reset the probability ratio)\;
			$g(\bs{r}',\bs{r},\Delta t)\leftarrow g(\bs{r},\bs{r}',\Delta t)$ (Reset the Green's function ratio) \;
		}
		{
			keep going\;
		}
	}
	\KwResult{The optimized trial wave function.}
	\caption{The Metropolis-Hastings algorithm. The positions are initialized randomly or was chosen by a previous sampling. The parameters are also usually initialized randomly or chosen by a parameter update. The Green's function ratio, $g$, can be evaluated efficiently using equation \eqref{eq:greensratio}, and the probability ratio $p$ can also often be found in a simple closed-form expression. The diffusion constant is $D=1/2$ in natural units. $\bs{\chi}$ is a random Gaussian variable. For more information, see section \ref{sec:importancesampling}}.
	\label{alg:hastings}
\end{algorithm}\DecMargin{1em}
\fi

\begin{lstlisting}[language=c++,caption={Taken from \lstinline{importancesampling.cpp}.}]
double ImportanceSampling::GreenRatio(const int i)
{
	double dQF = m_quantumForceOld(i) - m_quantumForceNew(i);
	return exp(0.5 * dQF * m_dx) + 1;
}
\end{lstlisting}
where \lstinline|dQF| is the difference between the new and the old force and \lstinline|m_dx| is the distance particle \lstinline|i| is moved. 

\section{Update of parameters} \label{sec:update}
The parameter update is a central part of a VMC implementation, and a good VMC implementation requires a good optimization algorithm. Since the optimization functions are called outside the sampling, they are just called a fraction number of times, compared to the function called from the sampling. Therefore, we will not put too much effort in making them efficient, but they should still be thought-through. We will here discuss the gradient descent method with momentum and monotonic decaying step, and the ADAM optimizer. 

The \lstinline|Optimizer| class contains a pure virtual function \lstinline|updateParameters| which is thus forced to be included in the optimizer subclasses. This function returns the update of the new parameters, and is the function we will discuss in this section.

\subsection{Gradient descent}
Gradient descent is a simple optimization algorithm, and so is the implementation. Based on the theory presented in section \ref{sec:gd}, the implementation is really straight-forward and reads
\begin{lstlisting}[language=c++,caption={Taken from \lstinline{gradientdescent.cpp}.}]
Eigen::MatrixXd GradientDescent::updateParameters()
{
	m_step += 1;
	double monotonic = 1 / pow(m_step, m_monotonicExp);
	m_v = m_gamma * m_v + m_eta * Optimization::getEnergyGradient() * monotonic;
	return m_v;
}
\end{lstlisting}
where \lstinline|m_v| is the momentum vector and \lstinline|m_monotonicExp| describes how fast the rate should decrease. Further, \lstinline|m_gamma| is the momentum parameter defining the relative size of the momentum. The function \lstinline|Optimization::getEnergyGradient| returns a matrix with the gradients of the energy expectation value with respect to all the parameters, given in equation \eqref{eq:gradientenergy}.

\subsection{ADAM optimizer}
The ADAM optimizer implementation also is very straight-forwardly based on the algorithm given in section \ref{sec:adam}. The momentum vectors were actually implemented as matrices to match the dimensions of the parameter matrix. By matrix operations we could also have made the function really fast, but since that is not the aim here, we decided to keep the loops in order to make the code legible. The implementation looks like:
\begin{lstlisting}[language=c++,caption={Taken from \lstinline{adam.cpp}.}]
Eigen::MatrixXd ADAM::updateParameters()
{
	m_step += 1;
	m_g = Optimization::getEnergyGradient();
	m_m = m_beta1 * m_m + (1 - m_beta1) * m_g;
	m_v = m_beta2 * m_v + (1 - m_beta2) * m_g.cwiseAbs2();
	m_mHat = m_m / (1 - pow(m_beta1, m_step));
	m_vHat = m_v / (1 - pow(m_beta2, m_step));
	for (int i = 0; i < m_numberOfElements; i++) {
		for (int j = 0; j < m_maxParameters; j++) {
			m_theta(i, j) = m_eta * m_mHat(i, j) / (sqrt(m_vHat(i, j) + m_epsilon));
		}
	}
	return m_theta;
}
\end{lstlisting}

The parameter matrix, named \lstinline|m_parameters|, can then easily be updated by the code
\begin{lstlisting}[language=c++,caption={Taken from \lstinline{system.cpp}.}]
if (m_myRank == 0) {
	m_sampler->computeAverages();
	m_parameters -= m_optimization->updateParameters();
}
\end{lstlisting}
where \lstinline|m_optimization| is the specified optimizer and \lstinline|m_myrank| is the \textit{rank} of the process. Parallel processing is not discussed yet, but we will describe it briefly in the following section.

\section{Parallel processing}
The code was parallelized using MPI to make it possible to study large systems. This means that the code can run multiple parallel treads and in that manner utilize the processors. Most notably, this allows us to run on computer clusters which typically .. We will not explain how MPI works in detail, nor will we detail the implementation of MPI since the commands are distribution over the entire code. The thing we present, is a sketch of the idea behind the parallelization used for our particular code.

Actually, one of the things that makes VMC preferred over other many-body methods, is that the algorithm quite easy can be split into independent parts, which encourages parallelization. The entire sampling can actually be split into as many parallel processes as needed, such that the code can be run on an arbitrary number of CPUs. 

We typically distinguish between wall clock time, $t_{clock}$ and CPU time $t_{cpu}$ where the former is the time measured by a clock and the latter is the total computation time from all the CPUs. The speed-up will in general not be 100\%, i.e., $t_{clock}\neq t_{cpu}/N$ with $N$ as the number of processes, mainly because all of the samplings should have the same burn-in period as if we only run one process, but also because the code that is not part of the sampling cannot be parallelized and needs to be run on a CPU. The process that takes care of this part is the main process with rank 0.

\IncMargin{1em}
\begin{algorithm}
	\SetAlgoLined
	MPI\_Init() (Initialize MPI)\;
	\While{not converged}{
		$E_L=0$\;
		gradient$=0$\;
		Egradient$=0$\;
		\For{$i\leftarrow 1$ \KwTo $M$}{
			$E_L+=(\hat{\mathcal{H}}\Psi)/\Psi$\;
			gradient$+=\nabla_{\theta}\ln\Psi$\;
			Egradient$+=(\hat{\mathcal{H}}\Psi)/\Psi*\nabla_{\theta}\ln\Psi$\;
		}
		\If{something goes wrong}{
			MPI\_Abort() (Abort all processes)\;
		}
		MPI\_Barrier() (Align processes)\;
		MPI\_Reduce($E_L$, gradients, Egradients) (Collect cumulative values)\;
		\If{myrank==0}{
			$\overline{E}_L = E_L/M$\;
			$\overline{\text{gradient}}=\text{gradient}/M$\;
			$\overline{\text{Egradient}}=\text{gradient}/M$\;
			$G=2*(\overline{\text{Egradient}}-\overline{E}_L*\overline{\text{Egradient}})$\;
			$\theta-=\eta G$\;
		}
		MPI\_Bcast($\theta$) (Broadcast parameters)\;
	}
	MPI\_Finalize() (Finalize MPI)\;
	\KwResult{Optimal variational parameters $\theta$.}
	\caption{Sketch of the parallelization.}
	\label{alg:mpi}
\end{algorithm}\DecMargin{1em}

In algorithm \ref{alg:mpi}, we have sketch very roughly how the parallelization goes. We first run the entire sampling individually for all the $N$ processes. If something goes wrong, we call the \lstinline|MPI_Abort| function. To align the processes before we collect all the cumulative values, we use the function \lstinline|MPI_Barrier| and we use \lstinline|MPI_Reduce| for the actual collection. Thereafter, the average energies are calculated \textit{by the main process only}, and in the end the updated parameters are broadcast to all the other processes. Note that this is just a sketch where we avoid the arguments and the actual implementation of the MPI functions. This is naturally found in the code. 

\section{Random number generator}
In Monte Carlo integration, we are dependent on random numbers which we get from a random number generator (RNG). The RNG should have to main properties: It should give many independent uncorrelated random numbers and it should be fast. The former depends on the \textit{period} of the RNG, where a long period gives many independent numbers. 

In this work, we have used the mersenne twister random number generator, as it has a period of $2^{19937}-1$ which is known as the Mersenne prime. This is a incredible large number, and should be more than sufficient for our purpose. We use built-in package in C++, \lstinline|std::mt19937|, which also is quite fast.